\documentclass[10pt]{article}
\usepackage[margin=0.72in]{geometry}
\usepackage{booktabs}
\usepackage{longtable}
\usepackage{array}
\usepackage[hidelinks]{hyperref}
\usepackage{microtype}
\setlength{\emergencystretch}{2em}

\title{TailGuard: Tail-Predictable Host-Memory Pooling for Disaggregated LLM Serving}
\author{Advisor-authored storyline draft; implementation and experiments remain student-owned}
\date{September 2026}

\begin{document}
\maketitle

\begin{abstract}
Disaggregated LLM serving can enlarge its KV cache by pooling CPU memory on
both prefill and decode nodes. The extra capacity is useful only if remote
reads, refill, and migration do not turn shared RDMA into an uncontrolled tail
latency path. We ask whether object placement and transfer budgets can be
jointly controlled so that both host-memory pools extend useful KV residence
while request p99 remains predictable. PegaFlow already contains host and SSD
backing, a vLLM connector, cross-node transport, and exploratory serving
traces. Those traces are not yet paper-admissible because historical code and
runtime custody is incomplete; they establish neither an RDMA-tail pathology
nor a benefit. We therefore formulate an object-lifecycle-aware placement and
transfer mechanism, its strongest baselines, and a fail-closed evaluation
contract. Cache TTL is treated as an outcome under a fixed total byte budget,
not a configurable success metric. No new performance or NPU claim is made in
this draft.
\end{abstract}

\section{Problem and Representative Setting}

Paged KV management makes cache capacity and fragmentation first-class serving
concerns~\cite{kwon2023pagedattention}. Systems such as Mooncake go further by
using CPU memory, SSD, and NIC resources across disaggregated prefill and decode
clusters~\cite{qin2025mooncake}. This creates a stronger baseline than
decode-local caching: both sides' memory can contribute to a distributed pool.
SYMPHONY additionally moves KV through a disaggregated memory layer using
advisory requests and cooperative cache management~\cite{agarwal2026symphony}.
It is therefore a required end-to-end baseline, not merely related work.

The representative setting is long-context P/D serving in which a decode node
frequently consumes KV produced elsewhere. A remote hit can avoid recompute,
but its critical path includes metadata lookup, memory registration, queueing,
RDMA transfer, destination placement, and materialization. Concurrent refill
or migration flows can share the same links with deadline-sensitive reads.
Consequently, more retained KV may coexist with worse request p99.

Our research question is:

\begin{quote}
Can a disaggregated KV store jointly choose host-memory placement, remote-read
timing, and transfer budgets so that it uses CPU memory on both P and D nodes
without making RDMA and request tail latency unpredictable?
\end{quote}

The paper advances only if a matched real-runtime trace first shows the claimed
counterexample: extra pooled capacity improves residency or hits while
uncontrolled remote traffic causes a repeatable tail-latency penalty.

\section{Gap and Scope}

Capacity-oriented policies decide what to retain; transfer engines decide how
to move bytes; serving schedulers decide where requests run. Treating these
decisions independently assumes that a remote cache hit has a stable cost.
That assumption can fail when large background movements and deadline-critical
reads share registration resources, queues, links, or destination bandwidth.

The intended contribution is not another multi-NIC routing protocol and not a
new KV correctness protocol. PegaFlow owns object placement, read admission,
and transfer budgeting at the cache/runtime boundary. Relative to SYMPHONY, the
claimed gap is narrower: whether one request-level receipt can expose and bound
the transfer debt jointly created by remote hits and background movement across
both P- and D-node host pools. Multi-path selection, connection recovery, and
generation-safe object identity remain required
components or baselines, but are outside the paper's novelty claim.

\section{Seven-Question Research Contract}

\begin{longtable}{>{\raggedright\arraybackslash}p{0.05\linewidth}p{0.90\linewidth}}
\toprule
Q & Frozen answer \\
\midrule
1 & \textbf{Problem.} Jointly exploit P- and D-node CPU memory while bounding
the tail cost of remote KV reads and background movement. \\
2 & \textbf{Importance.} A capacity win that violates TTFT or TPOT p99 is not a
serving win; long-context requests make both saved compute and transferred
bytes large. \\
3 & \textbf{Gap.} Existing capacity, placement, and transport controls do not
necessarily expose one request-level budget that accounts for queueing,
registration, transfer, refill, and materialization. Mooncake is a strong
end-to-end baseline, not evidence that the problem is already solved. \\
4 & \textbf{Idea.} Give every transferable KV object a lifecycle receipt and a
deadline-aware transfer value. Admit remote reads and background movement under
separate bounded budgets; place replicas according to predicted saved compute,
reuse distance, bytes, and contention. \\
5 & \textbf{Feasibility.} PegaFlow already exposes host backing, topology,
leases, prefetch, read cache, RDMA, connector, and metrics seams. The first
scope is one P/D deployment and one transfer path. \\
6 & \textbf{Evaluation.} Freeze total CPU memory, workload, topology, runtime,
and request order; compare local-only, static pooled memory, NUMA-aware
placement, transfer shaping, Mooncake-compatible policy, and the joint design. \\
7 & \textbf{Takeaway.} Distributed cache capacity must be scheduled together
with transfer debt; the useful region is predictable from saved compute,
reuse distance, object bytes, and deadline slack. \\
\bottomrule
\end{longtable}

\section{Mechanism Hypothesis}

For object $i$, the controller estimates
\[
V_i = C_i^{\mathrm{saved}} -
(L_i^{\mathrm{lookup}}+L_i^{\mathrm{queue}}+L_i^{\mathrm{xfer}}+
L_i^{\mathrm{materialize}}),
\]
along with bytes, expected reuse, deadline slack, source/destination pressure,
and placement lifetime. It chooses among local retain, remote retain, replicate,
defer, evict, or recompute. Deadline-critical reads receive a bounded foreground
budget; refill, replication, and migration consume a separately capped
background budget. A decision is revoked when identity, generation, topology,
or measured queue state differs from its admission receipt.

This is not merely rate limiting. Placement changes future traffic, while
traffic debt changes whether a remote copy is valuable. The hypothesis predicts
wins when reuse is high, saved prefill is large, and remote queues are
controllable. It predicts losses for short prefixes, one-shot objects, saturated
links, inaccurate reuse estimates, or materialization-dominated paths.

\section{Current Evidence Ledger}

The repository preserves real-online trace families and a portable trace-audit
harness. Issue 18 records that most historical referenced commits were not
reachable from durable refs at review time. Those artifacts remain exploratory:
they may select variables and failure cases, but no numeric result enters this
paper. Simulation and host probes qualify interfaces only.

\begin{table}[h]
\centering
\caption{Evidence boundary before the new matched run.}
\begin{tabular}{p{0.27\linewidth}p{0.25\linewidth}p{0.39\linewidth}}
\toprule
Artifact & Class & Admissible use \\
\midrule
Historical serving traces & exploratory real-online & variable selection; no number \\
Trace preregistration & contract & matched-arm definition \\
Host/simulation tests & qualification & semantics only \\
P/D RDMA tail counterexample & absent & no pathology claim \\
Joint policy result & absent & no benefit claim \\
\bottomrule
\end{tabular}
\end{table}

\section{Matched Evaluation Contract}

\paragraph{Trace schema.}
Each object records producer, consumer, model/request/generation identity,
bytes, source and destination tier, creation, lookup, queue, transfer,
materialization, hit/miss, reuse distance, deadline, cancellation, and cleanup.
NIC/link and NUMA identities must be tied to the request receipt.

\paragraph{Arms.}
Under identical model, requests, P/D counts, total CPU-memory bytes, topology,
and network conditions, compare: decode-local only; static use of both host
pools; NUMA-aware static placement; bounded-concurrency or priority shaping;
a deployable Mooncake-style baseline; and joint placement plus transfer budget.

\paragraph{Metrics.}
Report at least ten independent starts and full distributions for RDMA queue
delay, flow concurrency, bytes, useful bandwidth, cache hit, effective
residence time, TTFT, TPOT, throughput, recompute, recovery, and memory overhead.
All arms require completion, output, object identity, generation, cleanup, and
capacity conservation oracles.

\paragraph{Success gate.}
The primary gate requires at least a 20\% reduction in RDMA queue-delay p99 and
at least a 10\% reduction in request TTFT p99 relative to the strongest
deployable baseline, at unchanged total CPU memory. For each ratio, a
request-level stratified bootstrap over independent runs must place the upper
endpoint of the 95\% confidence interval below $0.80$ and $0.90$,
respectively. Cache hit may fall by at most one absolute percentage point and
throughput by at most 5\%; their one-sided 95\% confidence bounds must remain
inside those margins. These statistics and strata are frozen from baseline-only
pilots before treatment is run. A 50\% increase in effective cache residence is
reported only as a secondary capacity benefit and cannot compensate for worse
tail latency.

\section{Stopping Rules and Claim Discipline}

Stop the joint mechanism if matched real traces do not show an uncontrolled
traffic tail, if static placement or simple shaping reaches the oracle, or if
remote reads cannot beat recompute in the tested regime. Close only that
mechanism and retain the correctness carrier.

The present Ascend code establishes a porting/adaptation seam, not an
architecture-native contribution. An Ascend claim requires documented runtime
or memory behavior, a matched counterexample on 910B2, a native response, and a
predictive boundary. Host replay does not establish RDMA or serving benefit.

\section{Conclusion}

PegaFlow's paper question is not whether unused host memory can hold more KV.
It is whether distributed capacity can be made useful without importing
unbounded transfer debt into the request critical path. The next decisive
artifact is therefore a custody-complete P/D trace that jointly exposes
residency and tail cost, followed by the frozen matched comparison above.

\bibliographystyle{plain}
\bibliography{references}
\end{document}
